\chapter{Borders, padding, and box decoration}
\label{ch:borders}

\begin{aside}
The most basic question of UI: ``draw a box''. The answer
unfolds into Unicode rules, theme palettes, joinery for adjacent
borders, and the ancient art of corner glyphs.
\end{aside}

\section{Border characters}

Unicode block U+2500 BOX DRAWING gives us 128 glyphs for lines,
corners, junctions, and thick variants. \maya{}'s built-in
border presets are:

\begin{itemize}[leftmargin=*]
  \item \textbf{single.} Single thin: \code{U+2500}, \code{U+2502},
    \code{U+250C}, etc.
  \item \textbf{thick.} Heavy lines: \code{U+2501}, \code{U+2503},
    \code{U+250F}, etc.
  \item \textbf{double.} Double lines: \code{U+2550},
    \code{U+2551}, \code{U+2554}, etc.
  \item \textbf{rounded.} Single with rounded corners:
    \code{U+256D}, \code{U+256E}, \code{U+256F}, \code{U+2570}.
  \item \textbf{ascii.} Pure ASCII: \code{+}, \code{-}, \code{|}.
\end{itemize}

The \code{ascii} preset is for environments without Unicode
support (legacy terminals, dumb pipes for logging) and as a
diagnostic.

\section{The preset struct}

A border preset is six characters:

\begin{lstlisting}
struct BorderChars {
    char32_t top, bottom, left, right;
    char32_t tl, tr, bl, br;        // corners
};

constexpr BorderChars rounded = {
    U'\u2500', U'\u2500',           // top, bottom
    U'\u2502', U'\u2502',           // left, right
    U'\u256D', U'\u256E',           // top-left, top-right
    U'\u2570', U'\u256F',           // bot-left, bot-right
};
\end{lstlisting}

Custom presets are literally just a struct. The build-the-cool-
border ritual is: pick the eight characters, name them.

\section{Border placement}

Layout reserves one cell on each side for the border before
distributing space to children. A border element with size $w
\times h$ leaves $(w-2) \times (h-2)$ for its content.

This costs cells; for narrow elements (status bars, single-line
inputs) the border may not be worth it. \maya{}'s convention is
that the border is optional and off by default.

\section{Padding}

Padding is space between border and content. It does not draw
anything --- it just reserves cells:

\begin{lstlisting}
auto card = text("hi")
          | padding(1, 2)   // 1 row, 2 columns
          | border;
\end{lstlisting}

The padding goes inside the border. Margin (which is outside)
also exists and works the same way but reserves cells outside
the border.

\section{Joinery}

Two adjacent boxes can share their border. The classic case is a
table: every cell has its own border, but the borders touch and
should join with cross or tee glyphs.

\begin{lstlisting}
| A | B |
+---+---+
| C | D |
\end{lstlisting}

The \code{+} at the centre is a four-way junction. Each box on
its own would draw a corner; the table layout overlaps the
corners and substitutes the junction.

\maya{} performs this substitution in a post-paint pass:
\begin{itemize}[leftmargin=*]
  \item For each cell containing a border glyph, look at its
    four neighbours.
  \item If the neighbours' glyphs imply additional incoming
    strokes, upgrade this cell's glyph.
\end{itemize}

The upgrade table is 256 entries: for each combination of N, S,
E, W strokes (each present/absent/thick) pick the glyph.

\section{Thickness mixing}

A thick border next to a thin border needs a glyph that
transitions. \maya{}'s table has these (\code{U+2541},
\code{U+2542}, etc.) but only for single-to-thick.
Single-to-double does not exist in Unicode, so \maya{} falls
back to picking one side.

\begin{warning}
Mixing border presets in adjacent elements often looks worse
than picking one. The Unicode coverage of mixed-style
junctions is incomplete and what does exist looks uneven on
many fonts.
\end{warning}

\section{Theme colours}

A border carries a style: foreground, background, attributes.
The default takes the parent's foreground:

\begin{lstlisting}
auto panel = column(...)
           | border
           | fg<theme::muted>;   // muted border
\end{lstlisting}

For panels that need emphasis on focus, conditionally swap:

\begin{lstlisting}
auto panel = column(...)
           | border
           | fg<dynamic>(focused ? theme::accent
                                  : theme::muted);
\end{lstlisting}

\section{Titles}

A common pattern: a panel with a title embedded in the top
border:

\begin{lstlisting}
+-- Header ----------+
| body              |
+-------------------+
\end{lstlisting}

\maya{} provides \code{border\_with\_title(view)} which paints
the title at column 2, overwriting the border with the title's
text. The title can also take a style.

\section{Funky borders}

You can use any characters. Hashes \code{\#}, dashes \code{-},
ASCII art bricks, full-block characters, dotted lines (\code{U+
2508}, \code{U+250A}). Some apps use lighter colour for a
``ghost'' border that feels less heavy than the standard
single.

\begin{tryit}
Build a border preset using only spaces with background colour
set --- a coloured frame with no glyphs. Useful for high-
contrast UIs.
\end{tryit}

\section{No-border focus indication}

Borders are a heavy way to show focus. Lighter options:

\begin{itemize}[leftmargin=*]
  \item Change background to a slight tint.
  \item Add a left-edge stripe (one column of solid colour).
  \item Bold the text.
  \item Underline the active line.
\end{itemize}

Use heavy borders sparingly --- one main panel, perhaps, not
every input.

\section{Border bugs}

The two most common bugs:

\begin{itemize}[leftmargin=*]
  \item \textbf{Off-by-one.} Forgetting that the border eats
    cells; content overflows.
  \item \textbf{Single-cell width.} A border around a
    one-column element vanishes (border + padding + content =
    minimum width 3).
\end{itemize}

\maya{}'s layout engine refuses to draw a border smaller than 3
columns or rows: it logs a warning and renders the contents
without the border instead.

\section{Recap}

\begin{recap}
\begin{itemize}[leftmargin=*]
  \item Borders are presets: six or eight Unicode glyphs.
  \item Layout reserves cells around the content.
  \item Padding (inside) and margin (outside) tune the space.
  \item Joinery upgrades corner glyphs to junctions.
  \item Titles overlay the top border at column 2.
  \item Heavy borders are not the only focus indicator.
\end{itemize}
\end{recap}

\section*{Workshop}

\begin{enumerate}[leftmargin=*]
  \item Implement a custom border preset using \code{+},
    \code{=}, \code{|}, and ASCII corners. Apply it to a
    panel.
  \item Build a 3x3 grid of bordered cells. Confirm that
    \maya{}'s joinery upgrades each interior intersection to
    a cross.
  \item Replace one cell's border preset with the thick one
    and observe the transitions.
\end{enumerate}
